% Source: https://artofproblemsolving.com/community/c3h2398301_i_need_some_help_in_this_solid_geo_counting_problem
% Stable identity: aops-topic-2398301-q01
Let the polyhedron have $t$ triangular and $p$ pentagonal faces, so $t+p=32$.

\textbf{Counting corners.} Each triangle has $3$ corners and each vertex lies on $T$ triangles, so $3t=VT$; similarly $5p=VP$.

\textbf{Counting edges.} Each edge lies on two faces, so $2E=3t+5p=V(T+P)$.

\textbf{Euler.} $V-E+F=2$ becomes $V-\dfrac{V(T+P)}{2}+32=2$, i.e.
\[V(T+P-2)=60.\]
Also $t+p=32$ gives $\dfrac{VT}{3}+\dfrac{VP}{5}=32$, i.e. $V(5T+3P)=480$. Dividing the two equations,
\[\dfrac{5T+3P}{T+P-2}=8\iff3T+5P=16.\]
In nonnegative integers the only solution is $T=P=2$ ($P=0,1,3$ give $3T=16,11,1$). Then $V=\dfrac{60}{2}=30$, $t=20$, $p=12$, $E=60$, and indeed $30-60+32=2$. (This is the icosidodecahedron: around each vertex two triangles and two pentagons alternate, using $2\cdot60^\circ+2\cdot108^\circ=336^\circ$, as in the sketch.)

\begin{asy}
usepackage("amsmath");
size(4.5cm);
defaultpen(fontsize(10pt));
pair V=(0,0);
path regpoly(pair O, real alpha, int n){
  pair P=O; path g=P; real d=alpha;
  for(int i=0;i<n-1;++i){ P=P+dir(d); g=g--P; d+=360/n; }
  return g--cycle; }
filldraw(regpoly(V,0,3),rgb(0.75,0.85,1),black);
filldraw(regpoly(V,60,5),rgb(1,0.85,0.6),black);
filldraw(regpoly(V,168,3),rgb(0.75,0.85,1),black);
filldraw(regpoly(V,228,5),rgb(1,0.85,0.6),black);
dot(V);
label("$T$",V+0.6*dir(30)); label("$P$",V+1.25*dir(114)); label("$T$",V+0.6*dir(198)); label("$P$",V+1.25*dir(282));
draw(arc(V,0.45,336,360),red); label("$24^\circ$",V+0.75*dir(348),red+fontsize(8pt));
\end{asy}

So $100P+10T+V=200+20+30=\boxed{250}$.
